\IfFileExists{stacks-project.cls}{%
\documentclass{stacks-project}
}{%
\documentclass{amsart}
}

% For dealing with references we use the comment environment
\usepackage{verbatim}
\newenvironment{reference}{\comment}{\endcomment}
%\newenvironment{reference}{}{}
\newenvironment{slogan}{\comment}{\endcomment}
\newenvironment{history}{\comment}{\endcomment}

% For commutative diagrams we use Xy-pic
\usepackage[all]{xy}

% We use 2cell for 2-commutative diagrams.
\xyoption{2cell}
\UseAllTwocells

% We use multicol for the list of chapters between chapters
\usepackage{multicol}

% This is generally recommended for better output
\usepackage{lmodern}
\usepackage[T1]{fontenc}

% For cross-file-references
\usepackage{xr-hyper}

% Package for hypertext links:
\usepackage{hyperref}

% For any local file, say "hello.tex" you want to link to please
% use \externaldocument[hello-]{hello}
\externaldocument[introduction-]{introduction}
\externaldocument[conventions-]{conventions}
\externaldocument[sets-]{sets}
\externaldocument[categories-]{categories}
\externaldocument[topology-]{topology}
\externaldocument[sheaves-]{sheaves}
\externaldocument[sites-]{sites}
\externaldocument[stacks-]{stacks}
\externaldocument[fields-]{fields}
\externaldocument[algebra-]{algebra}
\externaldocument[brauer-]{brauer}
\externaldocument[homology-]{homology}
\externaldocument[derived-]{derived}
\externaldocument[simplicial-]{simplicial}
\externaldocument[more-algebra-]{more-algebra}
\externaldocument[smoothing-]{smoothing}
\externaldocument[modules-]{modules}
\externaldocument[sites-modules-]{sites-modules}
\externaldocument[injectives-]{injectives}
\externaldocument[cohomology-]{cohomology}
\externaldocument[sites-cohomology-]{sites-cohomology}
\externaldocument[dga-]{dga}
\externaldocument[dpa-]{dpa}
\externaldocument[sdga-]{sdga}
\externaldocument[hypercovering-]{hypercovering}
\externaldocument[schemes-]{schemes}
\externaldocument[constructions-]{constructions}
\externaldocument[properties-]{properties}
\externaldocument[morphisms-]{morphisms}
\externaldocument[coherent-]{coherent}
\externaldocument[divisors-]{divisors}
\externaldocument[limits-]{limits}
\externaldocument[varieties-]{varieties}
\externaldocument[topologies-]{topologies}
\externaldocument[descent-]{descent}
\externaldocument[perfect-]{perfect}
\externaldocument[more-morphisms-]{more-morphisms}
\externaldocument[flat-]{flat}
\externaldocument[groupoids-]{groupoids}
\externaldocument[more-groupoids-]{more-groupoids}
\externaldocument[etale-]{etale}
\externaldocument[chow-]{chow}
\externaldocument[intersection-]{intersection}
\externaldocument[pic-]{pic}
\externaldocument[weil-]{weil}
\externaldocument[adequate-]{adequate}
\externaldocument[dualizing-]{dualizing}
\externaldocument[duality-]{duality}
\externaldocument[discriminant-]{discriminant}
\externaldocument[derham-]{derham}
\externaldocument[local-cohomology-]{local-cohomology}
\externaldocument[algebraization-]{algebraization}
\externaldocument[curves-]{curves}
\externaldocument[resolve-]{resolve}
\externaldocument[models-]{models}
\externaldocument[functors-]{functors}
\externaldocument[equiv-]{equiv}
\externaldocument[pione-]{pione}
\externaldocument[etale-cohomology-]{etale-cohomology}
\externaldocument[proetale-]{proetale}
\externaldocument[relative-cycles-]{relative-cycles}
\externaldocument[more-etale-]{more-etale}
\externaldocument[trace-]{trace}
\externaldocument[crystalline-]{crystalline}
\externaldocument[spaces-]{spaces}
\externaldocument[spaces-properties-]{spaces-properties}
\externaldocument[spaces-morphisms-]{spaces-morphisms}
\externaldocument[decent-spaces-]{decent-spaces}
\externaldocument[spaces-cohomology-]{spaces-cohomology}
\externaldocument[spaces-limits-]{spaces-limits}
\externaldocument[spaces-divisors-]{spaces-divisors}
\externaldocument[spaces-over-fields-]{spaces-over-fields}
\externaldocument[spaces-topologies-]{spaces-topologies}
\externaldocument[spaces-descent-]{spaces-descent}
\externaldocument[spaces-perfect-]{spaces-perfect}
\externaldocument[spaces-more-morphisms-]{spaces-more-morphisms}
\externaldocument[spaces-flat-]{spaces-flat}
\externaldocument[spaces-groupoids-]{spaces-groupoids}
\externaldocument[spaces-more-groupoids-]{spaces-more-groupoids}
\externaldocument[bootstrap-]{bootstrap}
\externaldocument[spaces-pushouts-]{spaces-pushouts}
\externaldocument[spaces-chow-]{spaces-chow}
\externaldocument[groupoids-quotients-]{groupoids-quotients}
\externaldocument[spaces-more-cohomology-]{spaces-more-cohomology}
\externaldocument[spaces-simplicial-]{spaces-simplicial}
\externaldocument[spaces-duality-]{spaces-duality}
\externaldocument[formal-spaces-]{formal-spaces}
\externaldocument[restricted-]{restricted}
\externaldocument[spaces-resolve-]{spaces-resolve}
\externaldocument[formal-defos-]{formal-defos}
\externaldocument[defos-]{defos}
\externaldocument[cotangent-]{cotangent}
\externaldocument[examples-defos-]{examples-defos}
\externaldocument[algebraic-]{algebraic}
\externaldocument[examples-stacks-]{examples-stacks}
\externaldocument[stacks-sheaves-]{stacks-sheaves}
\externaldocument[criteria-]{criteria}
\externaldocument[artin-]{artin}
\externaldocument[quot-]{quot}
\externaldocument[stacks-properties-]{stacks-properties}
\externaldocument[stacks-morphisms-]{stacks-morphisms}
\externaldocument[stacks-limits-]{stacks-limits}
\externaldocument[stacks-cohomology-]{stacks-cohomology}
\externaldocument[stacks-perfect-]{stacks-perfect}
\externaldocument[stacks-introduction-]{stacks-introduction}
\externaldocument[stacks-more-morphisms-]{stacks-more-morphisms}
\externaldocument[stacks-geometry-]{stacks-geometry}
\externaldocument[moduli-]{moduli}
\externaldocument[moduli-curves-]{moduli-curves}
\externaldocument[examples-]{examples}
\externaldocument[exercises-]{exercises}
\externaldocument[guide-]{guide}
\externaldocument[desirables-]{desirables}
\externaldocument[coding-]{coding}
\externaldocument[obsolete-]{obsolete}
\externaldocument[fdl-]{fdl}
\externaldocument[index-]{index}

% Theorem environments.
%
\theoremstyle{plain}
\newtheorem{theorem}[subsection]{Theorem}
\newtheorem{proposition}[subsection]{Proposition}
\newtheorem{lemma}[subsection]{Lemma}

\theoremstyle{definition}
\newtheorem{definition}[subsection]{Definition}
\newtheorem{example}[subsection]{Example}
\newtheorem{exercise}[subsection]{Exercise}
\newtheorem{situation}[subsection]{Situation}

\theoremstyle{remark}
\newtheorem{remark}[subsection]{Remark}
\newtheorem{remarks}[subsection]{Remarks}

\numberwithin{equation}{subsection}

% Macros
%
\def\lim{\mathop{\mathrm{lim}}\nolimits}
\def\colim{\mathop{\mathrm{colim}}\nolimits}
\def\Spec{\mathop{\mathrm{Spec}}}
\def\Hom{\mathop{\mathrm{Hom}}\nolimits}
\def\Ext{\mathop{\mathrm{Ext}}\nolimits}
\def\SheafHom{\mathop{\mathcal{H}\!\mathit{om}}\nolimits}
\def\SheafExt{\mathop{\mathcal{E}\!\mathit{xt}}\nolimits}
\def\Sch{\mathit{Sch}}
\def\Mor{\mathop{\mathrm{Mor}}\nolimits}
\def\Ob{\mathop{\mathrm{Ob}}\nolimits}
\def\Sh{\mathop{\mathit{Sh}}\nolimits}
\def\NL{\mathop{N\!L}\nolimits}
\def\CH{\mathop{\mathrm{CH}}\nolimits}
\def\proetale{{pro\text{-}\acute{e}tale}}
\def\etale{{\acute{e}tale}}
\def\QCoh{\mathit{QCoh}}
\def\Ker{\mathop{\mathrm{Ker}}}
\def\Im{\mathop{\mathrm{Im}}}
\def\Coker{\mathop{\mathrm{Coker}}}
\def\Coim{\mathop{\mathrm{Coim}}}

% Boxtimes
%
\DeclareMathSymbol{\boxtimes}{\mathbin}{AMSa}{"02}

%
% Macros for moduli stacks/spaces
%
\def\QCohstack{\mathcal{QC}\!\mathit{oh}}
\def\Cohstack{\mathcal{C}\!\mathit{oh}}
\def\Spacesstack{\mathcal{S}\!\mathit{paces}}
\def\Quotfunctor{\mathrm{Quot}}
\def\Hilbfunctor{\mathrm{Hilb}}
\def\Curvesstack{\mathcal{C}\!\mathit{urves}}
\def\Polarizedstack{\mathcal{P}\!\mathit{olarized}}
\def\Complexesstack{\mathcal{C}\!\mathit{omplexes}}
% \Pic is the operator that assigns to X its picard group, usage \Pic(X)
% \Picardstack_{X/B} denotes the Picard stack of X over B
% \Picardfunctor_{X/B} denotes the Picard functor of X over B
\def\Pic{\mathop{\mathrm{Pic}}\nolimits}
\def\Picardstack{\mathcal{P}\!\mathit{ic}}
\def\Picardfunctor{\mathrm{Pic}}
\def\Deformationcategory{\mathcal{D}\!\mathit{ef}}
\usepackage{cleveref}
\crefname{originallemma}{original lemma}{original lemmas}
\Crefname{originallemma}{Original lemma}{Original lemmas}
\makeatletter
\newlabel{sites-cohomology-lemma-cartesian-flat-quasi-coherent}{{0GZR}{source}{}{sourcepointer-1}{}}
\newlabel{sites-cohomology-lemma-cartesian-flat-quasi-coherent@cref}{{[originallemma][0][]0GZR}{[1][0][]source}{}{}{}}
\makeatother
\crefname{originalremark}{original remark}{original remarks}
\Crefname{originalremark}{Original remark}{Original remarks}
\makeatletter
\newlabel{descent-remark-Zariski-site-space}{{03FG}{source}{}{sourcepointer-2}{}}
\newlabel{descent-remark-Zariski-site-space@cref}{{[originalremark][0][]03FG}{[1][0][]source}{}{}{}}
\makeatother
\crefname{originalproposition}{original proposition}{original propositions}
\Crefname{originalproposition}{Original proposition}{Original propositions}
\makeatletter
\newlabel{descent-proposition-equivalence-quasi-coherent}{{03DX}{source}{}{sourcepointer-3}{}}
\newlabel{descent-proposition-equivalence-quasi-coherent@cref}{{[originalproposition][0][]03DX}{[1][0][]source}{}{}{}}
\makeatother
\crefname{originallemma}{original lemma}{original lemmas}
\Crefname{originallemma}{Original lemma}{Original lemmas}
\makeatletter
\newlabel{sites-modules-lemma-special-locally-free}{{03DM}{source}{}{sourcepointer-4}{}}
\newlabel{sites-modules-lemma-special-locally-free@cref}{{[originallemma][0][]03DM}{[1][0][]source}{}{}{}}
\makeatother
\crefname{originallemma}{original lemma}{original lemmas}
\Crefname{originallemma}{Original lemma}{Original lemmas}
\makeatletter
\newlabel{sites-cohomology-lemma-compare-topologies-derived-adequate-modules}{{07A8}{source}{}{sourcepointer-5}{}}
\newlabel{sites-cohomology-lemma-compare-topologies-derived-adequate-modules@cref}{{[originallemma][0][]07A8}{[1][0][]source}{}{}{}}
\makeatother
\crefname{originallemma}{original lemma}{original lemmas}
\Crefname{originallemma}{Original lemma}{Original lemmas}
\makeatletter
\newlabel{descent-lemma-quasi-coherent-alternative-small}{{0GZX}{source}{}{sourcepointer-6}{}}
\newlabel{descent-lemma-quasi-coherent-alternative-small@cref}{{[originallemma][0][]0GZX}{[1][0][]source}{}{}{}}
\makeatother
\crefname{originallemma}{original lemma}{original lemmas}
\Crefname{originallemma}{Original lemma}{Original lemmas}
\makeatletter
\newlabel{sites-modules-lemma-chaotic-quasi-coherent}{{0GZN}{source}{}{sourcepointer-7}{}}
\newlabel{sites-modules-lemma-chaotic-quasi-coherent@cref}{{[originallemma][0][]0GZN}{[1][0][]source}{}{}{}}
\makeatother
\crefname{originallemma}{original lemma}{original lemmas}
\Crefname{originallemma}{Original lemma}{Original lemmas}
\makeatletter
\newlabel{coherent-lemma-quasi-coherent-affine-cohomology-zero}{{01XB}{source}{}{sourcepointer-8}{}}
\newlabel{coherent-lemma-quasi-coherent-affine-cohomology-zero@cref}{{[originallemma][0][]01XB}{[1][0][]source}{}{}{}}
\makeatother
\crefname{originalsection}{original section}{original sections}
\Crefname{originalsection}{Original section}{Original sections}
\makeatletter
\newlabel{schemes-section-quasi-coherent}{{01LA}{source}{}{sourcepointer-9}{}}
\newlabel{schemes-section-quasi-coherent@cref}{{[originalsection][0][]01LA}{[1][0][]source}{}{}{}}
\makeatother
\title{The five small quasi-coherent derived categories of a quasi-compact quasi-separated scheme agree}
\author{The Stacks Project Authors}
\date{stacks-project:perfect.tex:local-snapshot-20261012; GNU FDL 1.2 or later}
\begin{document}
\maketitle
\setcounter{subsection}{1}
\noindent\textit{Editorial note.} All delivered mathematical text is the author's own, copied from the pinned Stacks Project chapter file perfect.tex; no statement or proof step was rewritten, reordered or removed. The counted span is the author environment and its own proof environment. References inside the span are handed to the assembler as navigation stubs that carry the canonical tag of the cited tag only; whether each of them is load-bearing is recorded on the stub itself.\par\medskip
\noindent\textit{Allowed standard prerequisites (editorial).}\begin{itemize}
\item The category of sheaves on the small Zariski site of quasi-coherent modules agrees with the quasi-coherent modules on the ringed space, and the categories of modules agree, so the two derived categories are equivalent. invoked at proof lines 11019-11025. Gives the equivalence between the first two categories of the statement.
\item The equivalence of ringed topoi between the site of affines and the Zariski site induces an equivalence of module categories that preserves the quasi-coherent subcategories because intrinsicness of that notion. invoked at proof lines 11026-11032. Gives the equivalence between the second and third categories.
\item The comparison lemma for adequate modules produces an equivalence of derived categories from a morphism of ringed sites once four conditions on a basis and a subcategory of presheaves are checked. invoked at proof lines 11033-11060. Supplies the equivalence between the last two categories once its hypotheses are verified.
\item Quasi-coherent modules on the affine site are exactly the presheaves compatible with localizing the ring of sections, and they are also exactly the quasi-coherent modules in the chaotic topology. used at proof lines 11045-11058. Identifies the two descriptions of the subcategory to which the comparison lemma is applied.
\end{itemize}\medskip
\ifdefined\hypertarget\hypertarget{sourcepointer-1}{}\fi
\noindent\textit{Source reference.} Original lemma 0GZR: sources/stacks-chow/tags.txt, source lines 1--1. Cited at proof lines 11015-11017 for the equivalence between the last two categories; the flatness of the ring maps that makes it applicable is verified inside the counted span.\par\medskip
\ifdefined\hypertarget\hypertarget{sourcepointer-2}{}\fi
\noindent\textit{Source reference.} Original remark 03FG: sources/stacks-chow/tags.txt, source lines 1--1. Cited at proof lines 11019-11021 for the identification of the module categories; the subsequent correspondence of quasi-coherent subcategories is inside the counted span.\par\medskip
\ifdefined\hypertarget\hypertarget{sourcepointer-3}{}\fi
\noindent\textit{Source reference.} Original proposition 03DX: sources/stacks-chow/tags.txt, source lines 1--1. Cited at proof lines 11022-11024 to conclude the equivalence of the first two categories; everything else used for it is inside the counted span.\par\medskip
\ifdefined\hypertarget\hypertarget{sourcepointer-4}{}\fi
\noindent\textit{Source reference.} Original lemma 03DM: sources/stacks-chow/tags.txt, source lines 1--1. Cited at proof lines 11029-11031 to see that the equivalence of module categories preserves quasi-coherent subcategories; the conclusion drawn is inside the counted span.\par\medskip
\ifdefined\hypertarget\hypertarget{sourcepointer-5}{}\fi
\noindent\textit{Source reference.} Original lemma 07A8: sources/stacks-chow/tags.txt, source lines 1--1. Applied at proof lines 11033-11060 to produce the equivalence between the last two categories; the choice of basis and subcategory and the verification of the four conditions are inside the counted span.\par\medskip
\ifdefined\hypertarget\hypertarget{sourcepointer-6}{}\fi
\noindent\textit{Source reference.} Original lemma 0GZX: sources/stacks-chow/tags.txt, source lines 1--1. Cited at proof lines 11046-11048 to identify the members of the chosen subcategory with the quasi-coherent sheaves on the Zariski topology; the identification with the chaotic topology case follows inside the counted span.\par\medskip
\ifdefined\hypertarget\hypertarget{sourcepointer-7}{}\fi
\noindent\textit{Source reference.} Original lemma 0GZN: sources/stacks-chow/tags.txt, source lines 1--1. Cited at proof lines 11049-11053 for the second description of the subcategory; the conclusion that both descriptions agree is inside the counted span.\par\medskip
\ifdefined\hypertarget\hypertarget{sourcepointer-8}{}\fi
\noindent\textit{Source reference.} Original lemma 01XB: sources/stacks-chow/tags.txt, source lines 1--1. Cited at proof lines 11084-11086 for the fourth condition of the comparison lemma; the other three conditions are verified inside the counted span.\par\medskip
\ifdefined\hypertarget\hypertarget{sourcepointer-9}{}\fi
\noindent\textit{Source reference.} Original section 01LA: sources/stacks-chow/tags.txt, source lines 1--1. Cited at proof lines 11072-11076 for the second condition of the comparison lemma; the reduction to that discussion is inside the counted span.\par\medskip
\begin{lemma}
\label{lemma-DQCoh-alternative-small}
In the situation above there are canonical exact equivalences between
the following triangulated categories
\begin{enumerate}
\item $D_\QCoh(\mathcal{O}_X)$,
\item $D_\QCoh(X_{Zar}, \mathcal{O})$,
\item $D_\QCoh(X_{affine, Zar}, \mathcal{O})$,
\item $D_\QCoh(X_{affine}, \mathcal{O}_X)$, and
\item $\mathit{QC}(X_{affine}, \mathcal{O})$.
\end{enumerate}
\end{lemma}

\begin{proof}
If $U \subset V \subset X$ are affine open, then the ring map
$\mathcal{O}(V) \to \mathcal{O}(U)$ is flat. Hence
the equivalence between (4) and (5) is a special case of
Cohomology on Sites, Lemma
\ref{sites-cohomology-lemma-cartesian-flat-quasi-coherent}
(the proof also clarifies the statement).

\medskip\noindent
The ringed site $(X_{Zar}, \mathcal{O})$ and the ringed space
$(X, \mathcal{O}_X)$ have the same categories of modules by
Descent, Remark \ref{descent-remark-Zariski-site-space}.
Via this equivalence the quasi-coherent modules correspond by
Descent, Proposition \ref{descent-proposition-equivalence-quasi-coherent}.
Hence we get a canonical exact equivalence between the triangulated
categories in (1) and (2).

\medskip\noindent
The discussion preceding the lemma shows that we have an equivalence of
ringed topoi
$(\Sh(X_{affine, Zar}), \mathcal{O}) \to (\Sh(X_{Zar}), \mathcal{O})$
and hence an equivalence between abelian categories of modules.
Since the notion of quasi-coherent modules is intrinsic
(Modules on Sites, Lemma \ref{sites-modules-lemma-special-locally-free})
we see that this equivalence preserves the subcategories
of quasi-coherent modules. Thus  we get a canonical exact equivalence
between the triangulated categories in (2) and (3).

\medskip\noindent
To get an exact equivalence between the triangulated categories
in (3) and (4) we will apply Cohomology on Sites, Lemma
\ref{sites-cohomology-lemma-compare-topologies-derived-adequate-modules}
to the morphism $\epsilon : (X_{affine, Zar}, \mathcal{O}) \to
(X_{affine}, \mathcal{O})$ above. We take
$\mathcal{B} = \Ob(X_{affine})$ and we take
$\mathcal{A} \subset \textit{PMod}(X_{affine}, \mathcal{O})$
to be the full subcategory of those presheaves $\mathcal{F}$
such that $\mathcal{F}(V) \otimes_{\mathcal{O}(V)} \mathcal{O}(U)
\to \mathcal{F}(U)$ is an isomorphism. Observe that by
Descent, Lemma \ref{descent-lemma-quasi-coherent-alternative-small}
objects of $\mathcal{A}$ are exactly those sheaves in the Zariski
topology which are quasi-coherent modules on $(X_{affine, Zar}, \mathcal{O})$.
On the other hand, by Modules on Sites, Lemma
\ref{sites-modules-lemma-chaotic-quasi-coherent},
the objects of $\mathcal{A}$ are exactly the quasi-coherent modules
on $(X_{affine}, \mathcal{O})$, i.e., in the chaotic topology.
Thus if we show that Cohomology on Sites, Lemma
\ref{sites-cohomology-lemma-compare-topologies-derived-adequate-modules}
applies, then we do indeed get the canonical equivalence between the
categories of (3) and (4) using $\epsilon^*$ and $R\epsilon_*$.

\medskip\noindent
We have to verify 4 conditions:
\begin{enumerate}
\item Every object of $\mathcal{A}$ is a sheaf for the Zariski topology.
This we have seen above.
\item $\mathcal{A}$ is a weak Serre subcategory of
$\textit{Mod}(X_{affine, Zar}, \mathcal{O})$. Above we have
seen that $\mathcal{A} = \QCoh(X_{affine, Zar}, \mathcal{O})$ and
we have seen above that these,
via the equivalence $\textit{Mod}(X_{affine, Zar}, \mathcal{O}) =
\textit{Mod}(X, \mathcal{O}_X)$, correspond to the quasi-coherent modules
on $X$. Thus the result by the discussion in
Schemes, Section \ref{schemes-section-quasi-coherent}.
\item Every object of $X_{affine}$ has a covering in the chaotic
topology whose members are elements of $\mathcal{B}$. This holds
because $\mathcal{B}$ contains all objects.
\item For every object $U$ of $X_{affine}$ and $\mathcal{F}$ in $\mathcal{A}$
we have $H^p_{Zar}(U, \mathcal{F}) = 0$ for $p > 0$.
This holds by the vanishing of cohomology of quasi-coherent modules
on affines, see
Cohomology of Schemes, Lemma
\ref{coherent-lemma-quasi-coherent-affine-cohomology-zero}.
\end{enumerate}
This finishes the proof.
\end{proof}

\end{document}
